%% kapa: exam=1 round=27 year=2016 period=1 form=A subject=real_estate subject_ko=부동산학원론 seq=14 number=94
%% source: Q-Net 감정평가사 기출문제 (cst00309) artlSeq=5169402 file=extracted/1008/5169402/제27회 감정평가사 1차 문제/제27회 감정평가사 1교시문제(A형).pdf
%% conversion: kordoc 4.24.51 -> tools/gen1.py (auto); review=proofread-sonnet; flags=none
\begin{question}{94}
외부효과에 관한 설명으로 옳지 않은 것은?
\begin{choices}
  \item 외부효과는 한 사람의 행위가 제3자의 경제적 후생에 영향을 미치고, 그에 대해 지급된 보상을 제3자가 인지하지 못하는 현상을 말한다.
  \item 정(+)의 외부효과는 핌피(PIMFY) 현상을 초래할 수 있다.
  \item 부(-)의 외부효과를 완화하기 위한 수단으로 배출권 거래제도 등이 있다.
  \item 정(+)의 외부효과를 장려하기 위한 수단으로 보조금 지급 등이 있다.
  \item 공장이 설립된 인근지역에는 당해 공장에서 배출되는 폐수 등으로 인해 부(-)의 외부효과가 발생할 수 있다.
\end{choices}
\end{question}
\answer{1}
